\documentclass[11pt,a4paper,slovak]{article}

\usepackage[slovak]{babel}
\usepackage[T1]{fontenc}
\usepackage[utf8]{inputenc}

\usepackage[sfdefault]{FiraSans}

\usepackage{geometry}
\geometry{
    margin=2.5cm
}

\pagestyle{empty}
\title{Zameranie projektu}

\author{
{\small Metódy inžinierskej práce}\\
\\
{\small Fakulta informatiky a informačných technológií}\\
{\small Slovenská technická univerzita v Bratislave}\\
\\
{\small Bublincová Martina}\\
{\small Danečko Ján}\\
{\small Dobrev Andrej Thomas}\\
{\small Furman Michal}
}

\date{\small 4. október 2026}

\begin{document}
\pagestyle{empty}

\maketitle
\thispagestyle{empty}


\newpage


\section{Názov projektu}

Personalizovaný asistent sústredenia študentov



\section{Akronym} 
PASS 
\\
\\
''Use PASS to pass a test."

\section{Anotácia} 
Projekt PASS sa zameriava na problém digitálneho vyrušovania pozornosti študentov v dobe učenia. Mobilné telefóny sú častým distraktorom nielen kvôli notifikáciám, ale aj vďaka spontánnemu otváraniu aplikácií.  

Hlavným cieľom projektu je navrhnúť koncept mobilnej aplikácie určenej na zber, spracovanie a vizualizáciu údajov o rušení pozornosti počas učenia. Na základe identifikácie individuálnych vzorcov správania bude možné navrhnúť používateľovi vhodné metódy a  cielené opatrenia pre zníženie vplyvu rušivých digitálnych stimulov a návykov.  

Projekt bude zahŕňať  analýzu existujúcich riešení a súčasnej situácie výskumu v oblasti digitálneho rozptyľovania, návrh na efektívne vyhodnocovanie zozbieraných údajov a ich vzťahu s reakciami používateľa. Na rozdiel od bežných nástrojov založených na plošnom blokovaní notifikácií bude aplikácia PASS založená na individuálnom správaní a reakcii na konkrétne podnety, pričom bude používateľ informovaný o podvedomej manipulácii s mobilným telefónom a vyhľadávaní rozptýlenia.   

Navrhnutá aplikácia bude mať jednoduché používateľské rozhranie, s dvomi základnými maskami: Profil a Štúdium. Funkcia Štúdium aktivuje režim zberu dát a zapne reštrikcie podľa nastavení v sekcii Profil. Aplikácia zozbiera informácie o využívaní iných aplikácií, notifikáciách a reakciách používateľa. Po ukončení režimu Štúdium aplikácia uloží nazbierané údaje, spracuje ich a vyhodnotené informácie vráti do  sekcie Profil, kde ich používateľ môže sledovať a rozhodovať o nasledujúch opatreniach. Medzi ďalšie funkcionality budú patriť aj dlhodobé výsledky formou grafov, kalendár štúdia, body za nepretržitú sériu študijných relácií a časovače. 

Výsledkom bude koncept mobilnej aplikácie na systematickú podporu sústredenia počas štúdia, pričom rozhodnutia o konkrétnych krokoch a obmedzeniach aplikácií a notifikácií ostane vždy konečmým rozhodnutím používateľa.   

\nocite{*}
\bibliography{references}
\bibliographystyle{plain} 
\end{document}
